\documentclass[11pt]{article}
\usepackage[T1]{fontenc}
\usepackage[utf8]{inputenc}
\usepackage{lmodern}
\usepackage{amsmath,amssymb,amsthm,mathtools}
\usepackage{geometry}
\usepackage{hyperref}
\usepackage{enumitem}
\usepackage{bm}
\geometry{margin=1in}
\hypersetup{colorlinks=true,linkcolor=blue,citecolor=blue,urlcolor=blue}

\title{Actualized Geometry and a Null Prediction for the Bose--Marletto--Vedral Experiment in Relational Actualism}
\author{Joshua F. Sandeman}
\date{April 2026}

\newtheorem{definition}{Definition}
\newtheorem{lemma}{Lemma}
\newtheorem{theorem}{Theorem}
\newtheorem{corollary}{Corollary}
\newtheorem{remark}{Remark}
\newtheorem{proposition}{Proposition}

\begin{document}
\maketitle

\begin{abstract}
We show that Relational Actualism (RA) predicts a strict null result for gravity-mediated entanglement in the Bose--Marletto--Vedral (BMV) protocol. The reason is structural. In RA, the effective spacetime metric is a hydrodynamic description recalculated only from the actualized causal ledger, whereas the spatial branches of the BMV setup remain in the reversible possibility layer until positional actualization occurs. We formulate this Step-2 / Step-4 / Step-5 separation precisely and prove two results. First, in the branch-coherent regime no branch-resolved metric exists. Second, the gravitational phase acquired during the coherent interval is therefore common-mode across the four positional branches, so the BMV entangling invariant vanishes. We then clarify what ``spacetime superposition'' can and cannot mean in RA. The framework permits superposition over future ledger-growth channels that would yield different metrics upon actualization, but not superposition of the presently actual metric itself. A positive clean BMV signal would therefore directly falsify the present RA architecture.
\end{abstract}

\section{Introduction}

The Bose--Marletto--Vedral protocol has become a focal point in contemporary discussions of quantum gravity because it appears to offer a near-term experimental probe of a foundational question: can a quantum superposition of matter source a branch-dependent gravitational field? In the standard reading of the protocol, two mesoscopic masses are each prepared in a spatial superposition, and the gravitational interaction generates branch-dependent phases. If those phases are not additively separable, the two masses become entangled. On this reading, observation of the predicted entanglement is evidence that gravity itself must possess quantum degrees of freedom, or at least that gravity couples to quantum branch structure before any collapse-like event.

Relational Actualism makes a sharply different prediction. In RA, the metric is not a fundamental field defined over all possibilities, nor is it sourced by expectation values of unactualized branch structure. Rather, it is the macroscopic, hydrodynamic description of the already actualized causal Directed Acyclic Graph (DAG). Reversible quantum possibility and irreversible actualization are not two descriptions of the same ontic layer; they are distinct levels of the Engine of Becoming. Candidate futures are explored in a reversible amplitude layer, but the metric is recalculated only after an actualization event has been written into the ledger.

This distinction yields a direct consequence for the BMV protocol. The standard BMV entangling phase requires branch-dependent geometry during the interval in which the positional superpositions remain coherent. RA denies that such geometry exists. If the positions have not actualized, then the distinct branches are not distinct source distributions in the ledger; and if the positions do actualize, the coherence required for the interferometric protocol is destroyed. In RA, branch coherence and branch-resolved geometry are mutually exclusive in the BMV setting.

The argument below is intentionally narrow. We do not attempt a full review of BMV, nor do we discuss every variant of semiclassical or quantum-gravitational response. Our aim is more precise: to state the RA architecture cleanly, prove the null result at the level of a minimal theorem, and explain exactly what the framework means by rejecting ``spacetime superposition.''

\section{The standard BMV phase mechanism}

Consider two equal masses $A$ and $B$, each of mass $m$, prepared in the spatial superposition
\begin{equation}
|+\rangle_A = \frac{1}{\sqrt{2}}\bigl(|L\rangle_A + |R\rangle_A\bigr),
\qquad
|+\rangle_B = \frac{1}{\sqrt{2}}\bigl(|L\rangle_B + |R\rangle_B\bigr).
\end{equation}
The joint initial state is
\begin{equation}
|\Psi_0\rangle
=
|+\rangle_A \otimes |+\rangle_B
=
\frac{1}{2}\bigl(|LL\rangle + |LR\rangle + |RL\rangle + |RR\rangle\bigr).
\end{equation}

In the standard BMV account, each branch $|ij\rangle$ has separation $d_{ij}$ and therefore Newtonian interaction energy
\begin{equation}
U_{ij} = -\frac{Gm^2}{d_{ij}}.
\end{equation}
After interaction time $T$, the branch acquires phase
\begin{equation}
\phi_{ij} = -\frac{1}{\hbar} U_{ij} T = \frac{Gm^2 T}{\hbar\, d_{ij}}.
\end{equation}
Hence the evolved state is
\begin{equation}
|\Psi_{\mathrm{std}}(T)\rangle
=
\frac{1}{2}\sum_{i,j\in\{L,R\}} e^{i\phi_{ij}} |ij\rangle.
\end{equation}
The state is entangled when the phase is not additively separable:
\begin{equation}
\Delta\phi
:=
\phi_{LL} + \phi_{RR} - \phi_{LR} - \phi_{RL}
\neq 0 \pmod{2\pi}.
\end{equation}

This mechanism assumes that the four branch labels $LL,LR,RL,RR$ already correspond to four physically operative gravitational configurations during coherent evolution. The RA claim is that this assumption is false.

\section{The RA architecture: Step 2, Step 4, Step 5}

The relevant architectural distinction in RA can be stated succinctly.

\begin{itemize}[leftmargin=2em]
    \item \textbf{Step 2:} the reversible possibility layer, where amplitudes are assigned to candidate continuations of the current ledger.
    \item \textbf{Step 4:} actualization, in which one candidate event is irreversibly selected and written into the causal DAG.
    \item \textbf{Step 5:} the hydrodynamic limit, in which the effective metric is recalculated from the updated actualized graph.
\end{itemize}

The decisive claim is that Step 5 responds only to Step-4 actualizations, not to the Step-2 distribution of amplitudes across unactualized alternatives.

\begin{definition}[Actualized causal ledger]
Let $G_n = (V_n,E_n)$ denote the actualized causal DAG after $n$ actualization events. The vertex set $V_n$ contains only irreversible, ledger-writing actualization events; the edge set $E_n$ records their directed causal relations.
\end{definition}

\begin{definition}[Ledger-sourced metric]
The effective metric at stage $n$ is a functional of the actualized ledger:
\begin{equation}
 g^{(n)}_{\mu\nu} = \mathcal{G}[G_n],
\end{equation}
where $\mathcal{G}$ is the Step-5 coarse-graining map from the actualized graph to the emergent spacetime description.
\end{definition}

\begin{definition}[Branch-coherent BMV regime]
A BMV protocol is in the \emph{branch-coherent regime} on an interval $[0,T]$ if the positional alternatives $L$ and $R$ for each mass remain coherent throughout that interval, i.e. no irreversible positional actualization event occurs during $[0,T]$.
\end{definition}

Equivalently, the positional degree of freedom remains in the Step-2 possibility layer rather than passing through Step 4.

\section{The non-sourcing lemma}

We now formalize the core RA claim.

\begin{lemma}[Step-2 non-sourcing lemma]
Let $C_n$ denote the set of candidate extensions available from ledger state $G_n$ in the Step-2 possibility layer. Suppose a degree of freedom remains fully within the Step-2 layer during the interval $[0,T]$, so that none of its alternatives are selected at Step 4 during that interval. Then the Step-5 metric is insensitive to the decomposition of amplitudes across those alternatives.

Equivalently, if no candidate in the relevant sector actualizes on $[0,T]$, then
\begin{equation}
G(t) = G_n \qquad \text{for all } t\in[0,T],
\end{equation}
and therefore
\begin{equation}
g_{\mu\nu}(t) = \mathcal{G}[G_n] \qquad \text{for all } t\in[0,T].
\end{equation}
\end{lemma}

\begin{proof}
Step 2 enumerates weighted possibilities but does not itself add vertices to the actualized graph. By hypothesis, none of the relevant alternatives are selected at Step 4 during $[0,T]$. Hence the actualized graph in the relevant sector does not change on that interval. Since the Step-5 metric is a functional only of the actualized graph, it is unchanged throughout the interval. \qedhere
\end{proof}

\begin{corollary}[No branch-indexed metric before actualization]
In the branch-coherent BMV regime, there is no ontically meaningful family
\begin{equation}
\{g^{LL}_{\mu\nu}, g^{LR}_{\mu\nu}, g^{RL}_{\mu\nu}, g^{RR}_{\mu\nu}\}
\end{equation}
of simultaneously available branch-resolved metrics.
\end{corollary}

\begin{proof}
Such a family would require the labels $LL,LR,RL,RR$ to correspond to distinct actual source configurations already present in the ledger. But in the branch-coherent regime these labels denote only unactualized Step-2 alternatives. By the lemma, Step 5 responds only to the common actual ledger, not to branch labels that have not passed through Step 4. \qedhere
\end{proof}

\section{The RA-BMV null theorem}

We can now state the central result.

\begin{theorem}[RA-BMV null theorem]
Consider a BMV protocol with two masses $A$ and $B$, each prepared in a coherent two-position superposition. Assume:
\begin{enumerate}[label=(\arabic*),leftmargin=2.5em]
    \item the protocol remains in the branch-coherent regime during the interaction interval $[0,T]$;
    \item the effective metric is sourced only by the actualized ledger;
    \item nongravitational entangling channels are absent.
\end{enumerate}
Then the purely gravitational contribution to the four branch phases is common-mode:
\begin{equation}
\phi^{\mathrm{grav}}_{LL}
=
\phi^{\mathrm{grav}}_{LR}
=
\phi^{\mathrm{grav}}_{RL}
=
\phi^{\mathrm{grav}}_{RR},
\end{equation}
and therefore the entangling invariant vanishes:
\begin{equation}
\Delta\phi_{\mathrm{grav}}
:=
\phi^{\mathrm{grav}}_{LL}
+
\phi^{\mathrm{grav}}_{RR}
-
\phi^{\mathrm{grav}}_{LR}
-
\phi^{\mathrm{grav}}_{RL}
=
0.
\end{equation}
Hence the protocol cannot generate gravity-mediated entanglement.
\end{theorem}

\begin{proof}
By the corollary above, there is no branch-resolved metric during the coherent interval. Therefore the gravitational sector contributes the same Step-5 background to each amplitude component. The induced gravitational phase is identical on all four branches, so it factors as a global phase:
\begin{equation}
|\Psi_{\mathrm{RA}}(T)\rangle
=
 e^{i\phi_0}
\frac{1}{2}\bigl(|LL\rangle + |LR\rangle + |RL\rangle + |RR\rangle\bigr)
=
 e^{i\phi_0}|+\rangle_A\otimes |+\rangle_B.
\end{equation}
A global phase cannot entangle the state. Therefore the gravitational entangling invariant vanishes. \qedhere
\end{proof}

\begin{corollary}[Strict null prediction for ideal BMV]
In the idealized limit in which all nongravitational interactions are removed, the BMV entanglement witness must be null in RA.
\end{corollary}

\section{Coherence and geometry are mutually exclusive}

The theorem above rests on a deeper incompatibility.

\begin{theorem}[Coherence-geometry incompatibility theorem]
Within RA, the following two conditions cannot both hold for the same positional degree of freedom over the same interval:
\begin{enumerate}[label=(\alph*),leftmargin=2.5em]
    \item branch coherence: the alternatives $L$ and $R$ remain in coherent superposition;
    \item branch-resolved geometry: the alternatives $L$ and $R$ source distinct effective metrics.
\end{enumerate}
\end{theorem}

\begin{proof}
If branch coherence holds, then the alternatives remain in Step 2 and have not entered the ledger; by the non-sourcing lemma they cannot source distinct Step-5 metrics. Conversely, if branch-resolved geometry exists, then a branch label must already be represented in the actualized ledger, meaning Step 4 has occurred and the positional coherence has been destroyed. The two conditions therefore exclude one another. \qedhere
\end{proof}

This theorem is stronger than the slogan ``gravity is classical.'' It says, more precisely, that RA forbids simultaneous access to both the coherent superposition and the branch-dependent geometry required by the standard BMV mechanism.

\section{What `spacetime superposition' can and cannot mean in RA}

The phrase ``spacetime superposition'' is often used loosely. In the present context it must be split into two distinct meanings.

The \emph{forbidden} reading is:
\begin{quote}
The presently actual effective metric is itself in a coherent superposition of distinct branch geometries.
\end{quote}
This would amount schematically to a state of the form
\begin{equation}
\frac{1}{2}\Bigl(
|LL\rangle |g_{LL}\rangle
+
|LR\rangle |g_{LR}\rangle
+
|RL\rangle |g_{RL}\rangle
+
|RR\rangle |g_{RR}\rangle
\Bigr).
\end{equation}
RA denies that such a state is ontically meaningful, because the four metrics would require four distinct actual source configurations already present in the ledger.

The \emph{allowed} reading is weaker:
\begin{quote}
There is a coherent superposition of future ledger-growth channels, such that different eventual actualizations would yield different emergent metrics.
\end{quote}
In this sense RA permits superposition \emph{over geometry-producing futures}, but not superposition \emph{of presently actual geometry}. During the branch-coherent interval, the correct RA schematic is
\begin{equation}
\left(
\frac{1}{2}\sum_{i,j\in\{L,R\}} |ij\rangle
\right)
\otimes
|g[G_n]\rangle.
\end{equation}
Only the matter amplitudes are branch-resolved; the effective metric remains tied to the common actual ledger.

\section{Relation to semiclassical gravity}

The RA claim is not reproduced by naive semiclassical gravity. A semiclassical sourcing rule of the form
\begin{equation}
G_{\mu\nu} = 8\pi G\, \langle \hat{T}_{\mu\nu} \rangle
\end{equation}
allows unactualized branch structure to influence the field through expectation values. RA rejects this move at the ontological level. The source is not the expectation value of possibility; it is the actualized ledger of irreversible events.

Thus the RA null result should not be summarized merely as ``classical gravity.'' It is more specific: it is a consequence of \emph{actualized geometry}. Gravity in RA is the macroscopic description of what has already become actual.

\section{Experimental significance and falsification}

The logic of the BMV test is especially clean in RA.

\begin{itemize}[leftmargin=2em]
    \item If an ideal BMV experiment yields a null result for gravity-mediated entanglement after removal of nongravitational channels, the result is consistent with RA.
    \item If an ideal BMV experiment yields a positive clean signal, then branch-resolved geometry must be physically operative before actualization. That would collapse the Step-2 / Step-5 separation and directly falsify the present RA architecture.
\end{itemize}

This is not a small quantitative correction. A positive clean BMV signal would not merely require parameter adjustment; it would require a revision of the framework's core claim that the metric responds only to actualized vertices.

\section{Discussion}

The BMV protocol isolates a single question with unusual sharpness: does gravity couple to what is merely possible, or only to what has happened? Relational Actualism answers that it couples only to what has happened. The metric is not an additional quantum degree of freedom carried by every amplitude branch. It is the hydrodynamic record of the actualized causal DAG.

This distinction has two payoffs. First, it yields a clear null prediction. Second, it clarifies the conceptual status of ``spacetime superposition'' inside RA. The framework does not deny quantum superposition as such; it denies that the presently actual metric belongs to that superposition prior to actualization. Possibility superposes. The ledger does not.

For integration into the larger RA suite, the result is useful in two forms. In a trilogy or overview paper it can appear compactly as a theorem-level consequence of the Step-2 / Step-5 separation. As a stand-alone note, it has independent value because it directly engages an active experimental programme with a crisp, falsifiable prediction.

\section{Conclusion}

We have formulated the RA prediction for the Bose--Marletto--Vedral protocol in theorem form. The result is straightforward: in the branch-coherent regime, unactualized positional alternatives do not source branch-resolved geometry, because the effective metric is recalculated only from the actualized causal ledger. The gravitational phase is therefore common-mode across all four BMV branches, the entangling invariant vanishes, and gravity-mediated entanglement is forbidden. This null result is not incidental; it is a direct consequence of the Engine of Becoming's separation between reversible possibility and irreversible actualization. A positive clean BMV signal would therefore falsify the present RA architecture at its core.

\begin{thebibliography}{9}

\bibitem{Bose2017}
S. Bose, A. Mazumdar, G. Morley, H. Ulbricht, M. Toro\v{s}, M. Paternostro, A. Geraci, P. Barker, M. Kim, and G. Milburn,
\emph{Spin Entanglement Witness for Quantum Gravity},
Phys. Rev. Lett. \textbf{119}, 240401 (2017).

\bibitem{MarlettoVedral2017}
C. Marletto and V. Vedral,
\emph{Gravitationally Induced Entanglement between Two Massive Particles is Sufficient Evidence of Quantum Effects in Gravity},
Phys. Rev. Lett. \textbf{119}, 240402 (2017).

\bibitem{SandemanRA}
J. F. Sandeman,
\emph{Relational Actualism} suite of manuscripts, 2026.

\end{thebibliography}

\end{document}
